\scanpage{34}
dont le noyau \inserted{est $\bigcap_n nH^0(K,G)$, donc} est formé des éléments de $G(K)$ \uline{infiniment divisibles} dans $G$.

Or ici $K$ étant un \deleted{groupe} corps \inserted{absolument} de type fini, le th. de Mordell-Weil nous assure que $G(K)$ est \deleted{un groupe} \inserted{$\mathbb Z$-module de} type fini -- donc $G(K)\longrightarrow\varprojlim G(K)_n$ est injectif. Donc $G(K)\longrightarrow\varprojlim H^1(K,{}_nG)$ est injectif, et a fortiori l'hom (16), cqfd.

\uline{Remarque} Il ne suffirait pas dans cet argument de regarder, pour un $\ell$, l'extension de $E_{\overline K/K}$ par $H_1(X,\mathbb Z_\ell)$ -- il faut utiliser tous les $\ell$ -- si on omet $p$, alors la translation par un point d'ordre $p$ donne un contre-exemple. \deleted{\unsure{Donc}} Le résultat revient à ceci : Pour tt $x\in G(K)$, on regarde le torseur sur $K$, de groupe $T_\infty(G)=\varprojlim {}_nG$, image
\margin{33}
